\documentclass{article}
\usepackage{graphicx} % Required for inserting images
\usepackage{titlesec}
\usepackage[compress]{natbib}
\usepackage[resetlabels]{multibib}
\usepackage{authblk}
\usepackage{amsmath,amssymb,bm}
\usepackage{braket}

\usepackage[table]{xcolor}
\usepackage{url}            % simple URL typesetting
\usepackage{float}
\definecolor{dred}{rgb}{0.6,0,0}
\definecolor{dpurple}{HTML}{A020F0}
\definecolor{dblue}{rgb}{0,0,0.6}
\definecolor{hlcolor}{rgb}{1,1,0.8}
\usepackage[colorlinks=true, linkcolor=dblue, citecolor=dred, urlcolor=dblue]{hyperref}       % hyperlinks

\usepackage[nameinlink]{cleveref}
\usepackage[margin=1in]{geometry}
\renewcommand\b\bm
\newcommand*\sgn[1]{\text{sgn} ({#1})}
\setlength{\parindent}{0pt}
\setlength{\parskip  }{5.5pt}
\creflabelformat{equation}{#2\textup{#1}#3}

\title{\Large{Learning to predict model-based actions}}

\author[1]{\normalsize \bfseries Kristopher T.\ Jensen}
\affil[1]{ \small Sainsbury Wellcome Centre, University College London, London, UK}
\date{\vspace*{-2.0em}
\normalsize kris.torp.jensen@gmail.com
}



\begin{document}

\maketitle

Can we relax some of the assumptions about the planning process?

\subsection*{Task}

We consider the simple task of predicting the binary action of a linear teacher network from a $N$-dimensional input:
\begin{align}
    \b{x} &\sim \mathcal{N}(0, \b{I}_N)\\
    \b{w}^* &\sim \mathcal{N}(0, \b{I}_N / \sqrt{N})\\
    y^*(\b{x}) &:= \text{sgn}({\b{w}^*}^T \b{x}).
\end{align}

The student is similarly a linearly network with parameters $\b{w}$, and it produces a policy through a sigmoidal nonlinearity:
\begin{align}
    \pi(y = 1; \b{x}) &= \sigma(\b{w}^T \b{x})\\
    \pi(y=-1; \b{x}) &= 1-\sigma(\b{w}^T \b{x})\\
    \sigma(z) &= \frac{1}{1 + e^{-z}}.
\end{align}
A single trial consists of a sequence of $T$ independent `inputs' and corresponding actions.
The agent receives a reward of $R = 1$ only if the entire sequence of choices is correct:
\begin{equation}
    R(y_{1:T} ; y^*_{1:T}) = \prod_t [y_t = y^*_t].
\end{equation}

We denote the pre-activations $z = \b{w}^T \b{x}$ and $z^* = {\b{w}^*}^T \b{x}$.
Our pre-activations are jointly Gaussian with
\begin{align}
    \sigma_{|*}^2. := \braket{z z^*} &= \braket{\b{w}^T \b{x} \b{x}^T \b{w}^*} \\
    &= \b{w}^T \braket{\b{x} \b{x}^T } \b{w}^*\\
    &= \b{w}^T \b{w}^*\\
    \sigma_{||}^2 := \braket{z z} &= \b{w}^T \b{w} .\\
    \sigma_{**}^2 := \braket{z^* z^*} &= {\b{w}^*}^T \b{w}^*.
\end{align}
The correlation between $z$ and $z^*$ is given by $\rho_{|*} := \frac{\sigma_{|*}^2}{\sigma_{**} \sigma_{||}}$.
Without loss of generality, we will assume that the teacher is normalised such that $\sigma_{**}$ = 1$.

We can also calculate the probability that a greedy student action is correct, $p_g$, as a function of the student-teacher overlap $\rho$.
For a zero-mean bivariate Gaussian distribution, the orthant formula tells us that
\begin{equation}
    p(z >0, z^* > 0) = \frac{1}{4} + \frac{\text{arcsin}(\rho)}{2 \pi}.
\end{equation}
It follows that
\begin{align}
    p_g(\rho) &= p(\theta(z) = \theta(z^*))\\
    &= p(z >0, z^* > 0) + p(z < 0, z^* < 0)\\
    &= p(z >0, z^* > 0) + p(z >0, z^* > 0)\\
    &= \frac{1}{2} + \frac{\text{arcsin}(\rho)}{\pi}. \label{eq:pg}
\end{align}

We remind ourselves that the gradient of the logistic function is given by
\begin{align}
    \frac{d \sigma (z)}{dz} &= \frac{e^{-z}}{(1 + e^{-z})^2}\\
    &= \left ( \frac{1}{1 + e^{-z}} \right ) \left ( \frac{e^{-z}}{1 + e^{-z}} \right ) \\
    &= \sigma(z) (1 - \sigma(z)).
\end{align}
It follows that
\begin{align}
    \nabla_z \log \pi(y = 1) &= \nabla_z \log \sigma(z)\\
    &= (1 - \sigma (z))\\
    \nabla_z \log \pi (y = -1) &= \nabla_z \log (1 - \sigma(z))\\
    &= - \sigma (z)\\
    \nabla_z \log \pi (y = \tilde{y}) &= 0.5 + 0.5 \tilde{y} - \sigma(z).
    \label{eq:logsig_grad}
\end{align}
At various times, we will approximate the sigmoid of the final line as a step function ($\sigma(z) \approx \theta (z)$) or linearise it ($\sigma(z) \approx 0.5 + a z$) for analytical tractability.

We also note a few integrals that will prove useful later.
For $x \sim \mathcal{N}(0, 1)$:
\begin{align}
    \braket{\theta (\b{v}^T \b{x}) \b{x} } &= \frac{1}{2} \braket{ \hat{\b{v}}^T \b{x} | \b{v}^T \b{x} > 0 } \hat{\b{v}}\\
    &= \frac{1}{2} \braket{\text{abs}(z)}_{z \sim \mathcal{N}(0, 1)} \hat{\b{v}}\\
    &= \sqrt{\frac{1}{2 \pi}} \hat{\b{v}}\\
    &:= \beta \hat{\b{v}}.
\end{align}
Since $\text{sgn}(z) = 2 \theta(z) - 1$, it follows that $\braket{\text{sgn} (\b{v}^T \b{x}) \b{x} } = 2 \beta \tilde{\b{v}}$.
We can also evaluate similar integrals conditional on agreement between the teacher and a greedy student, with the student and teacher vectors separated by an angle $\theta = \text{arccos}(\rho)$:
\begin{align}
    \label{eq:conditonal_sgn}
    \braket{\text{sgn}(\b{v}^T \b{x}) \b{x} | \text{sgn}(\b{v}^T \b{x})= \text{sgn}(\b{u}^T \b{x})} &= \frac{\hat{\b{v}} + \hat{\b{u}}}{|| \hat{\b{v}} + \hat{\b{u}} || } \frac{\int_{r = 0}^{\infty} \int_{\phi = - \frac{\pi - \theta}{2}}^{\frac{\pi - \theta}{2}} e^{-r^2/2} \; r \cos \phi  \; r dr \; d \phi}{\int_{r = 0}^{\infty} \int_{\phi = - \frac{\pi - \theta}{2}}^{\frac{\pi - \theta}{2}} e^{-r^2/2} r dr d \phi} \\ 
    &= \frac{\hat{\b{v}} + \hat{\b{u}}}{ \sqrt{ 1 + 1 + 2 \cos \theta } } \frac{\sqrt{\pi / 2} [\sin \phi]_{\phi = - \frac{\pi - \theta}{2}}^{\frac{\pi - \theta}{2}} }{ [e^{-r^2/2}]_{r=0}^\infty [\phi]_{\phi = - \frac{\pi - \theta}{2}}^{\frac{\pi - \theta}{2}} }\\
    &= \frac{\hat{\b{v}} + \hat{\b{u}}}{ \sqrt{2 + 2 \rho} } \sqrt{\pi / 2} \frac{2 \sin (\frac{\pi - \theta}{2}) }{\pi - \theta}\\
    &= \frac{\hat{\b{v}} + \hat{\b{u}}}{ \sqrt{ 1 + \rho} } \sqrt{ \pi} \frac{\sin (\frac{\pi - \theta}{2}) }{\pi - \theta}\\
    &:= \kappa ( \hat{\b{v}} + \hat{\b{u}} ).
\end{align}
The first line follows from noting that $y = y^*$ condition is only satisfied for two cones with identical expectations, each of width $\pi - \theta$.
Within each cone, the vectors average out to the bisection, and the projection of a vector onto this line has magnitude $r \cos \phi$.
We integrate this projection over the area of interest.
Similarly,
\begin{align}
    \label{eq:conditonal_lin}
    \braket{(\hat{\b{v}}^T \b{x}) \b{x} | \text{sgn}(\b{v}^T \b{x})= \text{sgn}(\b{u}^T \b{x})} &= \hat{\b{v}} + \frac{\sin (\theta) }{\pi - \theta} \hat{\b{u}}\\
    &= \hat{\b{v}} + \frac{ \sqrt{1-\rho^2} }{\pi - \theta} \hat{\b{u}}.
\end{align}


\subsection*{Self-supervised learning}

We first consider the learning dynamics when doing supervised learning on training targets specified by an internal `planner' consisting of a simple sequence sampler and an exact evaluator.
We will learn from a target action $\tilde{y}$ via the following loss function:
\begin{equation}
    \mathcal{L}(\b{w}) := \mathbb{E}_{ \{ \b{x}_t \sim p(\b{x}_t) \} } \left [ T^{-1} \sum_t \log \pi_t (y_t = \tilde{y}_t) \right ].
\end{equation}
With a batch size of $L$ episodes, a weight update takes the form
\begin{align}
    \nabla_{\b{w}} \mathcal{L}(\b{w}) &= \frac{L T} \sum_{l=1, t=1}^{L, T} \nabla_{\b{w}} \log \pi_{lt} (y_{lt} = \tilde{y}_{lt}) \\
    &\approx \frac{L T} \sum_{l=1, t=1}^{L, T} \nabla_{\b{w}} (0.5-\theta(z_{lt})+0.5 \tilde{y}_{lt}) \b{x}_{lt}.
\end{align}
where we have approximate $\sigma(z) \approx [z > 0] := \theta(z)$.

I don't know to derive the learning dynamics unless I let $L \rightarrow$ infty, which corresponds to an infinite batch size (this may very well not be reasonable!).
In this case, the weight update is the mean of an infinite number of random variables, so it is just their expectation.
Given a probability $p_c$ of a single executed action being correct, we can write down this expectation:
\begin{align}
    \Delta \b{w} &= \braket{(0.5-\theta(z)+0.5 \tilde{y}) \b{x}}_{l, t} \\
    &= 0.5 \braket{\tilde{y} \b{x}} - \braket{\theta(z) \b{x}}\\
    &= 0.5 p_c \braket{y^* \b{x} | \tilde{y} = y^*} + 0.5 (1-p_c) \braket{ - y^* \b{x} | \tilde{y} \neq y^*} - \beta \hat{\b{w}}.
\end{align}
We will now make our first severe assumption, which is that $p(\b{x} | \tilde{y}=y^*) = p(\b{x}) = \mathcal{N}(0, \b{I}_N)$.
This implies that the errors of the planner are not confined to the same part of state space as the errors of the student!
In this case, we get
\begin{align}
    \Delta \b{w} & \approx 0.5 p_c \braket{y^* \b{x}} + 0.5 (p_c-1) \braket{ y^* \b{x}} - \beta \hat{\b{w}}\\
    &= (2 p_c - 1) \beta \hat{\b{w}}^* - \beta \hat{\b{w}}\\
    &= \alpha_{|} \hat{\b{w}} + \alpha_{*} \hat{\b{w}}^*,\\
    \alpha_{|} :&= - \beta\\
    \alpha_* :&= (2 p_c - 1) \beta.
\end{align}
The overlap with the student and target vectors are:
\begin{align}
    \label{eq:overlaps}
    (\Delta \b{w})^T \b{w} &= \alpha_| \sigma_{||} + \alpha_* \sigma_{|*}^2.\\
    (\Delta \b{w})^T \b{w}^* &= \alpha_| \sigma_{||}^{-1} \sigma_{|*}^2 + \alpha_*.\\
\end{align}

We now proceed to compute $p_c$.
To do so, we will make our second severe assumptionm which is that all simulated actions are sampled independently, and that they are correct with probability $p_gp = p(\theta(z) = \theta(z^*))$.
In other words, our sampler is \emph{as good as} the greedy student, but it can magically produce uncorrelated yet equally good samples.
This overestimates the prboability of sampling a correct sequence, especially when $p_g$ is at intermediate values.
In any case, we proceed undeterred.
The probability that at least one sampled sequence is entirely correct is:
\begin{equation}
    p_R := p( \exists_{i \in [1, 2, \ldots, M]} \; \text{s.t.} \; \tilde{y}_{it} = {y_t^*} \; \forall_t) = 1 - (1 - p_g^T)^M.
\end{equation}
In this case, all of the individual actions are correct.
The probability that no sampled sequences are entirely correct is $(1 - p_g^T)^M$.
In this case a random sequence sampled from the set of `not entirely correct' sequences is used as a target.
The marginal probability of an action being correct, given that the entire sequence is not, is:
\begin{align}
    p(y_t = y^* | \text{not all correct}) &= \frac{p(\text{not all correct}| y_t = y^*) p(y_t = y^*)}{p(\text{not all correct})}\\
    &= \frac{(1 - p_g^{T-1}) p_g}{1 - p_g^T}.
\end{align}
We therefore get
\begin{equation}
    \label{eq:pc}
    p_c = 1 - (1 - p_g^T)^M + (1 - p_g^T)^M \frac{(1 - p_g^{T-1}) p_g}{1 - p_g^T}.
\end{equation}

We now compute the change in $\sigma_{||}^2$ and $\sigma_{|*}^2$ when updating the weights with a learning rate $\eta$, $\b{w} \leftarrow \b{w} + \eta \Delta \b{w}$:

\begin{align}
    \label{eq:corr_change}
    \Delta \sigma_{|*}^2 &= (\b{w} + \eta \Delta \b{w})^T \b{w}^* - \b{w}^T \b{w}^* \\
    &= \eta (\Delta \b{w})^T \b{w}^* \\
    &= \eta \left [ \alpha_| \sigma_{||}^{-1} \sigma_{|*}^2 + \alpha_* \right ]\\
    &= \eta \beta \left [  (2 p_c - 1) - \rho \right ].
\end{align}

\begin{align}
    \label{eq:stud_change}
    \Delta \sigma_{||}^2 &= \Delta (\b{w}^T \b{w})\\
    &= (\b{w} + \eta \Delta \b{w})^T (\b{w} + \eta \Delta \b{w}) - \b{w}^T \b{w}\\
    &= 2 \eta (\Delta \b{w})^T \b{w} + \eta^2 (\Delta \b{w})^T (\Delta \b{w}) \\
    & 2 \eta \left [ \alpha_| \sigma_{||} + \alpha_* \sigma_{|*}^2 \right ] + \eta^2 \left [ \alpha_|^2 + \alpha_*^2 + 2 \alpha_| \alpha_* \rho \right ] \\
    &= 2 \eta \beta \left [  (2 p_c - 1) \sigma_{|*}^2 - \sigma_{||} \right ] + \eta^2 \beta^2 \left [  1  + (2p_c - 1)^2 - 2 (2 p_c - 1) \right ].
\end{align}
For small learning rates, we can ignore the last term, since $\eta^2 \rightarrow 0$.

Finally, we can write down the full learning dynamics:
\begin{align}
    \rho &= \frac{\sigma_{|*}^2}{\sigma_{||}}\\
    p_g &= \frac{1}{2} + \frac{\text{arcsin}(\rho)}{\pi}\\
    p_R &= 1 - (1 - p_g^T)^M\\
    p_c &= p_R + (1-p_R) \frac{(1 - p_g^{T-1}) p_g}{1 - p_g^T}\\
    \Delta \sigma_{|*}^2 &= \eta \beta \left [  (2 p_c - 1) - \rho \right ]\\
    \Delta \sigma_{||}^2 &= 2 \eta \beta \left [  (2 p_c - 1) \sigma_{|*}^2 - \sigma_{||} \right ] + \eta^2 \beta^2 \left [  1  + (2p_c - 1)^2 - 2 (2 p_c - 1) \right ].
\end{align}

\subsubsection*{Supervised learning}

We recover supervised learning simply by letting $M \rightarrow \infty$ so $p_c \rightarrow 1$.
In this case, we have
\begin{align}
    \Delta \sigma_{|*}^2 &= \eta \beta \left [  1 - \rho \right ]\\
    \Delta \sigma_{||}^2 &= 2 \eta \beta \left [  \sigma_{|*}^2 - \sigma_{||} \right ].
\end{align}

\subsubsection*{Reinforcement learning}

We now consider a reinforcement learner.
The policy gradient is given by:
\begin{equation}
    \nabla_{\b{w}} \mathcal{L}(\b{w}) = \mathbb{E}_{ \{ \b{x}_t \sim p(\b{x}_t) \} } \left [ T^{-1} \sum_t (R - B) \nabla_{\b{w}} \log \pi_t (y = y_t) \right ],
\end{equation}
where $B$ is a baseline.
For now, we will approximate $\sigma(z) = 0.5 + a \sigma_{||}^{-1} z := 0.5 + a \hat{z} $ in \Cref{eq:logsig_grad}.
With a batch size of $L$ episodes, a single gradient step takes the form
\begin{align}
    \Delta \b{w} &= \frac{L T} \sum_{l = 1}^L \sum_t (R_l - B) (0.5-\sigma(\hat{z}_{lt})+0.5 y_{lt}) \b{x}_{lt},\\
    \Delta \b{w} &\approx \frac{L T} \sum_{l = 1}^L \sum_t (R_l - B) (0.5 y_{lt} - a \hat{z}_{lt}) \b{x}_{lt}.
\end{align}
Nish's RL perceptron rule is recovered when $a = 0$.
We let $L \rightarrow$ infty, which corresponds to an infinite batch size (again this may very well not be reasonable!).
In this case, each gradient step is given by its expectation over time and batches:
\begin{align}
    \braket{\Delta \b{w}} &= \braket{(R - B) (0.5 y - a \hat{z}) \b{x}}_{l, t}\\
    &= 0.5 \braket{R y \b{x}} - a \braket{R \hat{z} \b{x}} + a B \braket{\hat{z} \b{x}} - 0.5 B \braket{y \b{x}}\\
    &= 0.5 p_R \braket{y \b{x} | y = y^*} - a p_R \braket{ \hat{z} \b{x} | y = y^*} + a B \braket{\hat{z} \b{x}} - 0.5 B \braket{y \b{x}}.
\end{align}
To evaluate the conditional expectations, we assume that the agent takes greedy actions, $y = \text{sgn}(z)$ (although we still assume a linearised sigmoid in the gradient, which is admittedly a bit weird).
Using \Cref{eq:conditonal_sgn} and \Cref{eq:conditonal_lin}, we get
\begin{align}
    \braket{\Delta \b{w}} &= 0.5 \kappa p_R (\hat{\b{w}} + \hat{\b{w}}^*) - a p_R \left (\hat{\b{w}} + \frac{\sqrt{1-\rho^2} }{\pi - \theta} \hat{\b{w}}^* \right ) + a B \hat{\b{w}} - B \beta \hat{\b{w}}\\
    &= \alpha_{|} \hat{\b{w}} + \alpha_{*} \hat{\b{w}}^*,\\
    \alpha_{|} :&= 0.5 p_R \kappa - a p_R  + a B - \beta B\\
    \alpha_* :&= 0.5 p_R \kappa - a p_R \frac{\sqrt{1-\rho^2} }{\pi - \theta}.
\end{align}

We can simply plug this result into \Cref{eq:overlaps}, \Cref{eq:corr_change}, and \Cref{eq:stud_change} to get:
\begin{align}
    \rho &= \frac{\sigma_{|*}^2}{\sigma_{||}}\\
    p_g &= \frac{1}{2} + \frac{\text{arcsin}(\rho)}{\pi}\\
    p_R &= p_g^T \\
    \Delta \sigma_{|*}^2 &= \eta \left ( \alpha_| \sigma_{||}^{-1} \sigma_{|*}^2 + \alpha_* \right ) \\
    \Delta \sigma_{||}^2 &= 2 \eta \left ( \alpha_| \sigma_{||} + \alpha_* \sigma_{|*}^2 \right ) + \eta^2 (\alpha_|^2 + \alpha_*^2 + 2 \alpha_| \alpha_* \rho).
\end{align}

If we let $a = 0$, we get
\begin{equation}
    \Delta \sigma_{|*}^2 &= 0.5 \eta p_g \kappa (1 + \frac{\sigma_{|*}^2}{\sigma_{||}^{-1} }) p_g^{T-1} - \eta \frac{B}{2 \pi} \frac{\sigma_{|*}^2}{\sigma_{||}^{-1} }.
\end{equation}
This recovers Nish's equation C25 up to a scaling (I don't know where I lost a factor of 2), with $r_1+r_2 = 1$, $r_2 = B$ and $p_g \kappa = \frac{1}{\sqrt{2 \pi}}$.
I didn't realise this last equality was true until I simulated it empirically, so clearly I have overcomplicated things somewhere but haven't looked into where yet.



\subsection*{Orthogonal learning}

In this section, we will assume normalised students and teachers and derive learning rules on the unit sphere.
We can write the \emph{orthogonal} component of the update as
\begin{align}
    \Delta \b{w}_o &= \Delta \b{w} - (\b{w}^T \Delta \b{w}) \b{w}.
\end{align}
Our expected updates all take the form
\begin{equation}
    \Delta \b{w} = \alpha_{|} \b{w} + \alpha_{*} \b{w}^*.
\end{equation}
The orthogonal component is
\begin{align}
    \Delta \b{w}_o &= \alpha_{|} \b{w} + \alpha_{*} \b{w}^* - (\alpha_{|} + \alpha_{*} \rho) \b{w}\\
    &= \alpha_{*} (\b{w}^* - \rho \b{w}).
\end{align}
The change in correlation is given by
\begin{align}
    \Delta \rho &= (\b{w} + \eta \Delta \b{w}_o)^T \b{w}^* - \b{w}^T \b{w}^*\\
    &= \eta \Delta \b{w}_o^T \b{w}^*\\
    &= (\alpha_{*} (\b{w}^* - \rho \b{w}))^T \b{w}^*\\
    &= \alpha_{*} (1 - \rho^2).
\end{align}

\subsubsection*{Gradient signal and noise}


Now let's compute the variance of our gradient estimators.
We will assume that $N \rightarrow \infty$, so $x_{it}$ is uncorrelated from $z$ and $y$, which neglects terms in $N^{-1}$.

\begin{align}
    \braket{\Delta w_i^2}_{x} &= \braket{((0.5 - \theta (z) + 0.5 \tilde{y})x_i)^2}\\
    &= 0.25 \braket{(\tilde{y} - y)x_{i}^2}\\
    &\approx 0.25 (\braket{\tilde{y}^2} + \braket{y^2} - 2 \braket{\tilde{y} y})\\
    &= 0.5 - 0.5 (p_c p_g + (1-p_c)(1-p_g) - p_g(1-p_c) - p_c(1-p_g))\\
    &= 0.5 - 0.5 (2 p_g - 1) (2 p_c - 1).
\end{align}
This holds for \emph{all} elements of the weight vector.
In the supervised learning case $p_c = 1.0$ and this reduces to $1 - p_g$.
At initialisation, $p_g = 0.5$ and we are left with $0.5$.
Importantly, the only effect of $T$ in the predictive learning case is to reduce the magnitude of the second term since $p_c \rightarrow 0.5$.
However, the second moment is still capped at 0.5.
The expected square of the mean updates is
\begin{align}
    \braket{\braket{\Delta w_i}_x^2}_i &= \braket{((2 p_c - 1) \beta \hat{w}^*_i - \beta \hat{w}_i)^2}_i\\
    &= N^{-1} \beta^2 ( 1 + (2p_c-1)^2 - 2 (2 p_c - 1) \rho ).
\end{align}
The $N^{-1}$ term arises because each element of our student and teacher scales with $N^{-\frac12}$ to ensure normalisation.
The scaled noise-to-signal ratio is
\begin{align}
    \gamma &= N^{-1} \frac{\braket{\Delta w_i^2}_{x} - \braket{\braket{\Delta w_i}_x^2}_i}{\braket{\braket{\Delta w_i}_x^2}_i}\\
    &\approx \frac{0.5 - 0.5 (2 p_g - 1) (2 p_c - 1)}{\beta^2 ( 1 + (2p_c-1)^2 - 2 (2 p_c - 1) \rho )}.
\end{align}
Notably, this is approximately constant in $T$.

In the RL case
\begin{align}
    \braket{\Delta w_i^2}_{x} &= \braket{((R-B)( 0.5 y - 0.25 \hat{z}) x_i)^2}\\
    &\approx \braket{(R-B)^2} (\frac{1}{4} \braket{y^2} + \frac{1}{16} \braket{\hat{z}^2} - \frac{1}{4} \braket{\text{sgn}(\hat{z}) \hat{z}})\\
    &= (p_R + B^2 - 2Bp_R)(\frac{5}{16} - \frac{1}{2}\beta)\\
    &\geq 0.113 p_R(1-p_R).
\end{align}
If we allow for a general linearisation $\sigma(z) \approx 0.5 + a z$, we get $p_R (1- p_R) (\frac{1}{4} + a^2 - 2 a \beta)$.
Conversely, 
\begin{align}
    \braket{\braket{\Delta w_i}_x^2}_i &= \braket{ \left ( \hat{\b{w}}_i \left ( 0.5 p_R \kappa - 0.25 p_R  + 0.25 B - \beta B \right ) + \hat{\b{w}^*}_i \left ( 0.5 p_R \kappa - 0.25 p_R \frac{\sqrt{1-\rho^2} }{\pi - \theta} \right ) \right )^2 }\\
    &= \left ( 0.5 p_R \kappa - \beta p_R \right )^2 + \left ( 0.5 p_R \kappa - 0.25 p_R \frac{\sqrt{1-\rho^2} }{\pi - \theta} \right )^2 + \rho \left ( 0.5 p_R \kappa - \beta p_R \right ) \left ( 0.5 p_R \kappa - 0.25 p_R \frac{\sqrt{1-\rho^2} }{\pi - \theta} \right )\\
    &= p_R^2 \left [ \left ( 0.5 \kappa - \beta \right )^2 + \left ( 0.5 \kappa - 0.25 \frac{\sqrt{1-\rho^2} }{\pi - \theta} \right )^2 + 2 \rho \left ( 0.5 \kappa - \beta \right ) \left ( 0.5 \kappa - 0.25 \frac{\sqrt{1-\rho^2} }{\pi - \theta} \right ) \right ].
\end{align}
Now the scaled signal-to-noise ratio is
The scaled noise-to-signal ratio is
\begin{align}
    \gamma &= N^{-1} \frac{\braket{\Delta w_i^2}_{x} - \braket{\braket{\Delta w_i}_x^2}_i}{\braket{\braket{\Delta w_i}_x^2}_i}\\
    &\approx p_R^{-1} \frac{0.11 (1-p_R)}{\left ( 0.5 \kappa - \beta \right )^2 + \left ( 0.5 \kappa - 0.25 \frac{\sqrt{1-\rho^2} }{\pi - \theta} \right )^2 + 2 \rho \left ( 0.5 \kappa - \beta \right ) \left ( 0.5 \kappa - 0.25 \frac{\sqrt{1-\rho^2} }{\pi - \theta} \right )}.
\end{align}
Notably, this scales as $p_R = p_g^T$!
We therefore see that both supervised and predictive learning have a constant signal-to-noise ratio, while RL is exponentially bad.


The gradient variance is exactly the same as in the non-orthogonal case.
This is intuitively because the gradients are order 1 times x, but the bit we subtract is multiplied by $\hat{w}$, which has elements of order $N^{-\frac12}$.
The gradient signal is given by
\begin{align}
    \braket{\braket{\Delta w_{oi}}_x^2 }_i &= \braket{ \left ( \alpha_{|} w_i + \alpha_{*} w^*_i - (\alpha_{|} + \alpha_{*} \rho) w_i \right )^2 }_i\\
    &= \braket{ \left ( \alpha_{|} w_i + \alpha_{*} w^*_i \right )^2  + \left ( (\alpha_{|} + \alpha_{*} \rho) w_i \right )^2 - 2 \left ( \alpha_{|} w_i + \alpha_{*} w^*_i \right ) (\alpha_{|} + \alpha_{*} \rho) w_i }_i\\
    &= \braket{ \alpha_{|}^2 w_i^2 + \alpha_{*}^2 w^*_i^2 + 2 \alpha_{|} \alpha_{*} w_i w^*_i + (\alpha_{|} + \alpha_{*} \rho)^2 w_i^2 - 2 \alpha_{|}^2 w_i^2 - 2 \alpha_{|} \alpha_{*} \rho w_i^2 - 2 \alpha_{|} \alpha_{*} w_i w_i^* - 2 \alpha_{*}^2 \rho w_i w_i^*}_i\\
    &= \alpha_{|}^2 + \alpha_{*}^2 + 2 \alpha_{|} \alpha_{*} \rho + \alpha_{|}^2 + \alpha_{*}^2 \rho^2 + 2 \alpha_{|} \alpha_{*} \rho - 2 \alpha_{|}^2 - 2 \alpha_{|} \alpha_{*} \rho - 2 \alpha_{|} \alpha_{*} \rho - 2 \alpha_{*}^2 \rho^2\\
    &= \alpha_{*}^2 + \alpha_{*}^2 \rho^2 - 2 \alpha_{*}^2 \rho^2\\
    &= \alpha_{*}^2(1 - \rho^2).
\end{align}




\end{document}